\documentclass[12pt]{article}
\usepackage[letterpaper,margin=1in]{geometry}
\usepackage{fontspec}
\usepackage{setspace}
\usepackage{amsmath}
\usepackage{amssymb}
\usepackage{microtype}
\usepackage{xcolor}
\usepackage{fancyhdr}
\usepackage{titlesec}
\usepackage{parskip}
\usepackage{hyperref}
\hypersetup{
  colorlinks=true,
  linkcolor=blue!50!black,
  urlcolor=blue!50!black,
  citecolor=blue!50!black,
  pdftitle={A1 to A2: The Actualization of Perception},
  pdfauthor={Dalton Salvo},
  pdfsubject={Dissertation -- Aristotelian Actualization Chain}
}

% Main font: Linux Libertine if available, otherwise default serif
\IfFontExistsTF{Linux Libertine O}{%
  \setmainfont{Linux Libertine O}[Numbers=OldStyle]%
}{%
  \IfFontExistsTF{TeX Gyre Pagella}{%
    \setmainfont{TeX Gyre Pagella}%
  }{}%
}

% Greek wrapping: try dedicated Greek font, else fall back to main font
\IfFontExistsTF{GFS Didot}{%
  \newfontfamily{\greekfont}{GFS Didot}%
  \newcommand{\gk}[1]{{\greekfont #1}}%
}{%
  \IfFontExistsTF{DejaVu Serif}{%
    \newfontfamily{\greekfont}{DejaVu Serif}%
    \newcommand{\gk}[1]{{\greekfont #1}}%
  }{%
    % Last resort: use main font (works if it has Greek glyphs)
    \newcommand{\gk}[1]{#1}%
  }
}

% Line spacing and paragraph style
\setstretch{1.3}
\setlength{\parindent}{0pt}
\setlength{\parskip}{0.8em}

% Section titles: generous spacing, clean typography
\titleformat{\section}{\Large\bfseries}{}{0pt}{}
\titleformat{\subsection}{\large\bfseries}{}{0pt}{}
\titlespacing*{\section}{0pt}{2em}{1em}
\titlespacing*{\subsection}{0pt}{1.5em}{0.6em}

% Footnote style: slightly tighter
\usepackage[hang,flushmargin]{footmisc}

% Page header/footer
\setlength{\headheight}{14pt}
\pagestyle{fancy}
\fancyhf{}
\fancyhead[L]{\small\itshape A\textsubscript{1}$\to$A\textsubscript{2}: Aisthēsis}
\fancyhead[R]{\small\itshape Draft — \today}
\fancyfoot[C]{\thepage}
\renewcommand{\headrulewidth}{0.4pt}
\renewcommand{\footrulewidth}{0pt}

% Title page content
\title{\vspace{-2em}\textbf{$A_1 \to A_2$: The Actualization of Perception}\\[0.5em]
       \large\normalfont A Dissertation Section (Console Draft)}
\author{Dalton Salvo}
\date{April 2026 — Draft for Review}

\begin{document}
\maketitle

\thispagestyle{fancy}

\subsection*{1. $A_1$→$A_2$: The Actualization of Perception}

$A_0$ established motion and time as the first actuality of the chain --- the ontological ground whose ongoing activity licenses every subsequent stage. The present section specifies how perception actualizes within that ground. Two preconditions together constitute $A_1$: (i) the sensible object actually exercises its sensible power --- the colour visible, the sound audible, the warm warming --- and (ii) the sense-organ holds its first actuality as the capacity to receive sensible form without matter. The motion that carries $A_1$ to its completion --- \textit{aisthēsis}, the perceptual act --- is $A_1$$\to$$A_2$. Its terminus, $A_2$, is the completed perceptual actuality: the joint actualization of sensible object and sense faculty in which the organ receives the form of its object without receiving the object's matter, together with the residual motion deposited in the sensory apparatus when the perceptual event ends.

The section proceeds through seven architectural moves. First, it specifies the $A_1$ preconditions and the medium that transmits the sensible form. Second, it develops Aristotle's hylomorphic account of form-reception through the wax-signet metaphor of \textit{De Anima} II.12. Third, it applies the three-factor kinetic schema established at $A_0$ to perception, showing that the activity of the sensible object and the activity of the sense faculty are numerically one event. Fourth, it analyses perception as an enmattered \textit{alloiōsis} of the preservative kind rather than the corrupting kind. Fifth, it establishes that the perceptual motion does not terminate with the object's withdrawal but persists as residual motion in the sensory apparatus --- what this chapter names resonant \textit{kinēsis}. Sixth, it advances the dual-trace thesis: the residual motion is not a single undifferentiated impression but a composite whose formal-epistemic aspect (resonant \textit{aisthēma}) and affective-orectic aspect (resonant \textit{orexis}) are one in substrate and two in being. Seventh, it situates $A_2$'s dual-trace composite as the unmoved originator of the phantastic motion $M_2$$\to$$A_3$. A Development Note flags three open questions that subsequent sections will take up.

The architectural stakes are considerable. If $A_2$'s completed actuality harboured only a formal trace, the chain could not explain how perception generates the motivational states that drive animal action. The dual-trace thesis developed in §7 is thus not an ornament but a structural requirement: it specifies the material from which the chain's subsequent stages --- the \textit{phantasma} at $A_3$, the cognitive actualities and the committal dimension of \textit{doxa} at $M_2$$\to$$M_3$, the determinate \textit{pathē} that mobilize \textit{orexis} at $M_3$$\to$$M_4$ --- will be generated. Aristotle does not name the second trace; the reading advanced here treats the dual structure as a structural requirement of his own commitments, to be defended by close engagement with the texts that ground it.

\subsection*{2. The Preconditions: $A_1$ and the Medium}

For $A_1$$\to$$A_2$ to occur, two actualities must already obtain. The first is the sensible quality exercising its power as sensible. A colour is visible when it is actually colouring --- when the surface is exhibiting its chromatic form in the presence of the medium that will transmit it. A sound is audible when it is actually sounding --- when the struck object or resonant body is exercising the capacity that its material structure enables. Aristotle does not develop a separate technical analysis of sensible qualities in act, because the point is already available from the $A_0$ framework: the sensible object is an \textit{entelecheia} whose \textit{telos} is internal, and its exercise requires no further enabling condition beyond the motion and time in which all natural activity occurs.

The second $A_1$ precondition is the sense-organ in its first actuality. Aristotle develops the first/second actuality distinction most clearly at \textit{De Anima} II.5, 417a21-b2. The geometer who holds geometrical knowledge without currently exercising it possesses a first actuality --- a \textit{hexis}, a stable disposition capable of being exercised. The geometer who actively proves a theorem exercises that capacity and thereby possesses a second actuality. The sense-organ's first actuality is likewise a \textit{hexis}: the eye's capacity to see, the ear's capacity to hear, the tongue's capacity to taste. This first actuality is not a bare potentiality to be awakened from without; it is already an actuality, a structured readiness that awaits only the proper object to pass into the second actuality of exercise. $A_1$, accordingly, is not a stage of pure potentiality but a stage at which two actualities --- sensible object and sense-organ --- both obtain but have not yet entered into the joint actualization that $A_1$$\to$$A_2$ will complete.

Between the sensible object and the sense-organ stands the medium --- the air, water, or transparent substance through which the sensible form is transmitted. Aristotle develops the medium's role at \textit{Sense and Sensibilia} 6, 447a3-7: "if the body which is heated or frozen is extensive, each part of it successively is affected by the part contiguous, while the part first changed in quality is so changed by the cause itself which originates the change." The medium's affection is sequential. The sensible quality alters the part of the medium in contact with the object; that altered part alters the next part; and so on until the alteration reaches the sense-organ. Sound, whose propagation Aristotle analyses most explicitly, is genuinely successive: we perceive the clap of thunder after we see the flash because the alteration of the air takes measurable time to traverse the medium. Light is the exception. Aristotle notes at \textit{Sense and Sensibilia} 6, 446b8-9 that light is transmitted instantaneously --- not propagated successively through the medium but actualizing the transparent in a single movement. The sequential character characteristic of sound and heat does not obtain for vision.

The medium's sequential affection is the physical correlate of the temporal horizon $A_0$ establishes. Where $A_0$ supplies the general ontological conditions --- motion and time as co-actual --- the medium supplies the specific physical mechanism through which one particular species of motion, the perceptual kind, unfolds in time. The sequential character matters architecturally because it means that even the first moment of perception is temporally extended. Reception is not an instantaneous event occurring at a mathematical point; it is a temporally unfolding process whose temporality is inherited from $A_0$. The three actualities of $A_1$ --- object, organ, medium --- are each possible only because motion and time are already in act. $A_1$, accordingly, is not an independent stage but a specification of $A_0$'s horizon as it obtains for living beings with sense-organs and sensible objects in the field of perception.

\subsection*{3. Hylomorphism and the Reception of Form}

The perceptual motion from $A_1$ to $A_2$ is a distinctive species of \textit{alloiōsis} in which the sense-organ receives the form of the sensible object without receiving the object's matter. The governing text is \textit{De Anima} II.12, 424a17-24: "a sense is what has the power of receiving into itself the sensible forms of things without the matter, in the way in which a piece of wax takes on the impress of a signet-ring without the iron or gold; what produces the impression is a signet of bronze or gold, but not \textit{qua} bronze or gold: in a similar way the sense is affected by what is coloured or flavoured or sounding not insofar as each is what it is, but insofar as it is of such and such a sort and according to its form." The signet-ring analogy is architecturally decisive. The ring impresses its engraved pattern into the wax; the wax takes the pattern; but the wax does not become iron or gold. What passes from ring to wax is the \textit{form} --- the structured pattern of convexities and concavities --- not the substrate that bears the form. Perception operates analogously. The colour of the apple passes to the eye; the pitch of the string passes to the ear; but the eye does not become apple-matter, nor the ear string-matter. The form is received as form, and its reception is the actualization of the sense-organ's capacity to receive that form.

Aristotle distinguishes the organ as spatial magnitude from the sensory power as form-in-magnitude at \textit{De Anima} II.12, 424a24-29: "The sense and its organ are the same in fact, but their essence is not the same. What perceives is, of course, a spatial magnitude, but we must not admit that either the having the power to perceive or the sense itself is a magnitude; what they are is a certain form or power in a magnitude." The distinction is crucial. The eye is a spatial magnitude --- a physical object with extension, composed of matter, located in the body. The \textit{capacity} to see, however, is not itself a magnitude; it is a form present in the magnitude of the eye, a ratio or structured configuration that the organ's material substrate realizes. When the organ receives the sensible form, it is the form-in-magnitude that is altered, not the material magnitude. The material magnitude is unchanged in kind; what changes is the formal configuration that the material magnitude realizes.

The form-in-magnitude that is the sensory power, Aristotle argues, is a \textit{logos} --- a ratio of the contraries proper to the sense's domain. At \textit{De Anima} II.11, 424a4-5 he observes that sense is "a sort of mean" between the extremes of the sensible qualities, and at \textit{De Anima} III.2, 426a27-b8 he develops the ratio analysis explicitly. Each sense is constituted by a proportion: the eye between light and dark, the ear between high and low pitch, the tongue between sweet and bitter. The \textit{logos} is structurally a ratio, and perceptible qualities are themselves ratios of their underlying contraries --- gray as a specific proportion of white and black, a flavour as a proportion of sweet and bitter. Perception, on this account, is the organ's taking-on of the ratio that the perceptible quality instantiates: the sensory power adjusts its ratio to match the object's ratio, and this formal matching \textit{is} the perceptual act.

The \textit{logos} framework explains the destruction condition that Aristotle develops at \textit{De Anima} III.2, 426b3-8: if the sensible exceeds the ratio of the organ, the motion acts on the organ \textit{as matter} rather than \textit{as sense}. Excess brightness destroys sight; excess loudness destroys hearing; excess heat destroys touch. When the stimulus intensity exceeds the proportional range the organ can receive, the organ is affected as matter rather than as sense, and the \textit{logos} itself is destroyed rather than realized. The motion produced is real, and it is genuine damage to the organ, but it is not perception. Perception requires that the stimulus fall within the ratio the organ can receive; excess falls outside the ratio, and the organ suffers rather than perceives. This destruction condition presupposes the hylomorphic analysis: if perception were simply material alteration, there would be no principled difference between perceiving and being damaged. The \textit{logos} is what makes that difference principled.

The wax-signet analogy accordingly does more than illustrate form-reception; it marks the specific mode of being that perception exhibits. The wax undergoes genuine alteration --- it acquires a pattern it did not previously have --- but its material substrate is preserved rather than transformed. The alteration is formal, not material. This points forward to the next architectural move. Perception is an \textit{alloiōsis}, but not every \textit{alloiōsis} is corruption; the distinction between two kinds of qualitative alteration, and the assignment of perception to the second kind, must now be developed.

\subsection*{4. The Activity of Mover and Moved as One}

The three-factor kinetic schema established at $A_0$ --- unmoved originator, moved mover, that which is moved --- applies directly to perception. The unmoved originator at $A_1$$\to$$A_2$ is the sensible object in actuality, whose form initiates the motion without being itself altered by it. The moved mover is the sense faculty, which is altered (formally, per §3) in the very act of receiving the sensible form and which thereby moves the perceiver as a whole. That which is moved is the perceiver via the sense-organ. The schema iterates from $A_0$ with the same structural logic: one actuality functions as the unmoved source of the motion that generates the next actuality.

Decisive for the perceptual application is Aristotle's claim at \textit{Physics} III.2, 202a13-20 that the actuality of the mover and the actuality of the moved are numerically one. "Motion is in the movable. It is the fulfilment of this potentiality by the action of that which has the power of causing motion; and the actuality of that which has the power of causing motion is not other than the actuality of the movable; for it must be the fulfilment of \textit{both}." One motion, not two. The mover's \textit{poiesis} and the moved's \textit{pathesis} are distinguished in account but not in substrate; they are one event described under two categorial aspects. Aristotle makes the point concrete with the Thebes-Athens analogy at \textit{Physics} III.3, 202b13-14: "the road from Thebes to Athens and the road from Athens to Thebes are the same." One road, two descriptions. One motion, two formal aspects.

\textit{De Anima} III.2, 425b26-27 extends this directly to perception: "the activity of the sensible object and that of the sense is one and the same activity, and yet the distinction between their being remains." The sensible object's exercise of its power (its \textit{poiesis}) and the sense faculty's reception of the form (its \textit{pathesis}) are not two events conjoined but one event under two \textit{logoi}. Aristotle illustrates with a striking example at \textit{De Anima} III.2, 425b28-426a1: "a man may have hearing and yet not be hearing, and that which has a sound is not always sounding. But when that which can hear is actively hearing and that which can sound is sounding, then the actual hearing and the actual sound come about at the same time." Hearkening and sounding --- the active hearing and the active sounding --- arise together, in a single actualization. The sound is not first produced and then heard; the sounding and the hearkening are the same actuality under two descriptions.

The formula Aristotle employs here --- one in substrate, different in being --- will recur at every subsequent stage of the chain. It is introduced at the motion-level in \textit{Physics} III.2-3 (mover and moved, \textit{poiesis} and \textit{pathesis}); it is redeployed for perception at \textit{De Anima} III.2 (sensible and sense); it will be applied again at \textit{De Anima} III.7, 431a13-14 to the coordination of perception and appetite (one in substrate, different in being); and it structures the dual-trace thesis to be developed in §6. The formula's persistence across contexts reflects the chain's architectural unity: a single ontological pattern --- numerical identity under plural formal descriptions --- iterates across levels, each time specifying a new relation that the substrate-level unity secures.

One further feature of the perceptual motion must be registered. Because the mover's actuality and the moved's actuality are numerically one, the perceptual event does not consist in the mover producing something that then exists separately from the mover in the moved. The actuality \textit{is} the shared event, and the moved's being-affected is an aspect of that event rather than a subsequent state it acquires. This structural feature licenses what §6 will develop: the event, once complete, does not simply vanish. Motion that is the shared actuality of two substances leaves traces in the moved that are the residual character of that shared actuality. The next two subsections trace what happens as the perceptual motion completes and what remains when the sensible object withdraws.

\subsection*{5. Perception as Enmattered Alloiōsis}

Aristotle frames perception as a species of \textit{alloiōsis} --- qualitative alteration --- at \textit{De Anima} II.5, 416b33-417a1: "Sensation depends, as we have said, on a process of movement or affection from without, for it is held to be some sort of change in quality." Two kinds of \textit{alloiōsis} must be distinguished. The first is ordinary qualitative change, in which a substance's quality is altered by the imposition of a contrary: cold water becomes hot, a living body becomes diseased, a flame is extinguished. This kind of \textit{alloiōsis} is corrupting in the sense that it destroys the subject's prior state; the contrary supplants what was there.

The second kind of \textit{alloiōsis} --- which Aristotle develops at \textit{De Anima} II.5, 417a21-b2 --- is preservative rather than corrupting. "When we use the term 'to be acted on,' we do not use it in one single sense, but in two: one is the destruction of one of two contraries by the other, the other is the preservation of what is potential by what is actual and like it." The geometer learning geometry, the builder acquiring the art of building, the sense-organ receiving a sensible form: these are alterations that do not destroy the subject's prior state but actualize a potentiality that was already present as first actuality. The preservative \textit{alloiōsis} is not a process in which the subject is replaced by its contrary; it is a process in which a capacity already present in the subject is brought to exercise.

Perception belongs to the preservative kind. When the eye sees red, the eye is not thereby destroyed or replaced by red; its capacity to see is actualized, and the actualization is a \textit{fulfilment} of that capacity rather than its negation. Aristotle is explicit at \textit{De Anima} III.7, 431a3-5: "in the case of sense clearly the sensitive faculty already was potentially what the object makes it to be actually; the faculty is not affected or altered" in the corruptive sense. The sensitive faculty receives the sensible form by actualizing what it was already capable of being; the reception does not introduce a foreign element that destroys the organ's prior state.

This preservative character is what makes perception a genuine actualization rather than a deficient alteration. The wax-signet analogy clarifies the point. The wax is altered by the ring --- it acquires a pattern it did not previously have --- but the alteration is not the wax's corruption; the wax is preserved as wax, and the pattern it receives is a determination of the wax rather than a substitution for it. The sensory organ is likewise preserved as what it is while receiving the form of what it perceives; the reception is a determination of the organ, not a substitution.

The structural consequence matters decisively for what follows. If perception were a corruptive \textit{alloiōsis} --- if the organ were destroyed or substantially transformed in each perceptual act --- then nothing would remain of the act once it ended. The organ would have to be restored before the next perception could occur. But perception is not corruptive; it is preservative, and the preservation means that the form received leaves a trace. The organ, having been formally altered, does not return to its pre-perceptual state the moment the object withdraws. Something persists. What persists, and how it does so, is the subject of the next subsection.

\subsection*{6. Resonant Kinesis: The Persisting Impression}

The perceptual motion does not terminate when the sensible object withdraws. Aristotle makes this point at several places in the corpus with striking consistency. At \textit{De Anima} III.2, 425b23-26 he observes: "even when the sensible objects are gone the sensings and imaginings continue to exist in the sense-organs." At \textit{On Dreams} 2, 459a24-28 he elaborates: "The objects of sense-perception corresponding to each sensory organ produce sense-perception in us, and the affection due to their operation is present in the organs of sense not only when the perceptions are actualized, but even when they have departed." The residual affection is not a metaphor for an after-effect; it is the continuation of the perceptual motion in the organ after the originating cause has withdrawn.

Aristotle illustrates the mechanism with a physical analogy at \textit{On Dreams} 2, 459b1-6: "This we must likewise assume to happen in the case of qualitative change; for that part which has been heated by something hot, heats the part next to it, and this propagates the affection onwards to the starting-point. This must therefore happen in sense-perception, since actual perceiving is a qualitative change." The analogy compares the residual motion to heat remaining in water after the fire is removed. The water, having been heated, does not instantaneously return to its pre-heated state when the fire is withdrawn; the warmth persists as a residual motion in the water's material substrate, gradually dissipating as the water equilibrates with its surroundings. So too in perception: the organ, having been formally altered, retains the alteration as a residual motion that persists for some time after the object has withdrawn.

Aristotle gives a name to these residual affections at \textit{On Dreams} 2, 460b2-3: "even when the external object of perception has departed, the impressions it has made persist, and are themselves objects of perception." The Greek is \textit{\gk{αἰσθήματα}} --- \textit{aisthēmata}, impressions or residual sensings. The \textit{aisthēmata} are the persisting formal traces of the completed perceptual act. They are "themselves objects of perception" in the sense that they can be attended to by the sensory apparatus in the original object's absence; they are the material of dreams, afterimages, imagination, and the beginnings of memory.

This chapter adopts the term *\textit{resonant }kinēsis**\textit{ for the residual motion understood as a whole --- the ongoing activity in the sensory apparatus that persists after the sensible object has withdrawn. The coinage marks two features of the trace that Aristotle's own lexicon does not consolidate into a single technical vocabulary. First, the trace is }kinēsis\textit{ --- it is a motion, not a static imprint, and it has the ongoing character of all Aristotelian }kinēsis\textit{. Second, it is }resonant\textit{ --- it continues reverberating within the organ like sound continuing in a stilled bell, retaining the formal structure of the original motion while operating in the absence of its originating cause. Kevin White captures the reverberating character precisely: "This lingering, resonating, echoing presence of sensible forms freed from their original matter appears to be the essence of }phantasia\textit{ for Aristotle. The movement in the sense-power caused by the sensible object itself sets up a second movement which continues after the sensation has ceased and the object is gone" (White, }The Meaning of Phantasia in Aristotle's De Anima III, 3-8*, p. 498). The residual motion is not a weakened copy of the original; it is a continuing exercise of the same formal structure, sustained by the organ's material substrate in the original cause's absence.

The signet-ring analogy, extended to three phases, illuminates what resonant \textit{kinēsis} adds to the hylomorphic account of §3. Phase one: potentiality. The wax is pliable and smooth; the ring is held above. The sense-organ is ready and the sensible object nearby but unseen --- $A_1$ is complete, $A_2$ not yet underway. Phase two: actualization. The ring presses into the wax; the ring and wax are jointly actual in the impressing; "the actual hearing and the actual sound come about at the same time," as Aristotle puts it at \textit{De Anima} III.2, 425b28-426a1. This is $A_2$ itself: the perceptual act as the shared actuality of sensible object and sense faculty. Phase three: residual trace. The ring is withdrawn, the actualization proper ends --- but the wax has undergone preservative \textit{alloiōsis}, and the impression persists. The \textit{aisthēmata} remain. The resonant \textit{kinēsis} is the organ's continuing exercise of the form received, sustained in the object's absence by the organ's own material substrate.

The structural role of resonant \textit{kinēsis} in the chain follows from its character. An alteration that vanished with its cause could not ground further actualities; a completed actuality that terminated in itself could not serve as the unmoved originator of the next motion. But resonant \textit{kinēsis} persists, and in persisting, it becomes the material upon which the phantastic motion $M_2$$\to$$A_3$ will operate. The \textit{phantasma} that will emerge at $A_3$ is not created ex nihilo from the sensible object's form; it is generated from the resonant \textit{kinēsis} that $A_2$ has deposited in the organ. $A_2$'s residual trace is thus architecturally indispensable: without it, the chain could not continue past the perceptual event. The character of that residue, however --- whether it is a single undifferentiated impression or a structured composite --- determines what the subsequent stages of the chain can generate from it. The next subsection argues that the residue is not single but dual.

\subsection*{7. The Dual-Trace Thesis: Resonant Aisthēma and Resonant Orexis}

Here the section does its most consequential philosophical work. The resonant \textit{kinēsis} that persists after the perceptual event is not a single undifferentiated impression but a composite with two formally distinct aspects. The first aspect --- the resonant \textit{aisthēma} --- is the formal-epistemic residue: the shape, colour, pitch, or scent of the perceived object preserved in the organ as the "what it is" of the original sensation. The second aspect --- the resonant \textit{orexis} --- is the affective-orectic residue: the pleasure or pain, the pursue or avoid, that accompanied the original perceptual act and persists alongside the formal trace as the "how it matters" for the organism's being. The dual trace is numerically one event in the organ's substrate and two in being --- the same formula Aristotle applies to motion at \textit{Physics} III.2, to perception at \textit{De Anima} III.2, and now, as this section argues, to the residual motion that perception leaves behind.

The argument proceeds along two converging lines. The first is analytic. Aristotle establishes at \textit{De Anima} II.3, 414b1-6 a chain of entailments: "whatever has a sense has the capacity for pleasure and pain and therefore has pleasant and painful objects present to it, and wherever these are present, there is desire, for desire is appetition of what is pleasant." Every act of sensation, in this chain, analytically entails affective valence. To sense at all is already to be oriented toward the pleasant and away from the painful; to have pleasant and painful objects present is to have the appetite those objects engage. Sensation, on this account, is not a neutral registration of qualities that can then be evaluated by some further faculty; sensation is already evaluative in its very structure.

Aristotle develops the coordination of perception and appetite directly at \textit{De Anima} III.7, 431a8-14: "To perceive then is like bare asserting or thinking; but when the object is pleasant or painful, the soul makes a sort of affirmation or negation, and pursues or avoids the object. To feel pleasure or pain is to act with the sensitive mean towards what is good or bad as such. Both avoidance and appetite when actual are identical with this: the faculty of appetite and avoidance are not different, either from one another or from the faculty of sense-perception; but their being is different." The passage is decisive. Perception and appetite are numerically one in substrate --- one faculty, one event --- but they differ in being. When the object is pleasant or painful, the perceptual act is simultaneously an affirmation or negation, a pursuit or avoidance. There is not first a perception and then, by a separate cognitive operation, an evaluation; the perception itself carries the evaluation as a formal aspect of its being.

The residual motion that persists after the perceptual act must inherit this dual structure. The sensation that deposited the trace was itself numerically one and two in being; its \textit{pathos} was formally both perceptual and orectic. Unless one can argue that the residue discards one aspect while retaining the other --- which Aristotle's hylomorphism does not permit --- the full \textit{pathos} persists, with both formal aspects intact. The enmattered-accounts argument at \textit{De Anima} I.1, 403a25-b19 supports this directly. The affections of soul are "enmattered accounts" (\textit{logoi enyloi}) whose definition must combine formal and material aspects; the physicist gives the material account (a boiling of the blood around the heart), the dialectician the formal account (a desire for retaliation), and the full definition requires both. If the residual motion is an enmattered \textit{pathos}, and the \textit{pathos} has two formally distinct aspects (perceptual and orectic), then the residue preserves both aspects as formal aspects of a single material substrate.

The second line of argument is a modus tollens. If the residual motion were exhausted by its formal-epistemic aspect --- if what persisted were only the resonant \textit{aisthēma} --- then two consequences would follow that Aristotle's texts decisively reject. First, mere retention of sensible form would suffice to originate pursuit and avoidance. The soul's knowledge \textit{that the object is a wolf} would, by itself, move the animal to flee. But Aristotle insists that bare perception is "like bare asserting or thinking"; it is only \textit{when the object is pleasant or painful} that the soul "makes a sort of affirmation or negation, and pursues or avoids the object" (\textit{De Anima} III.7, 431a8-14). The formal content alone is motivationally inert. A distinct operation --- the hedonic-orectic valuation of the object as pleasant or painful --- is required to initiate pursuit or avoidance. If the residue carried only the formal trace, it could not motivate; something more must persist.

Second, if only the formal trace persisted, the affective residue could have no independent operative force. But Aristotle observes at \textit{On Dreams} 2, 460b3-11 that "when we are excited by emotions" we are deceived by slight formal resemblances: "the coward when excited by fear, the amorous person by the object of his desire; so that, with but little resemblance to go upon, the former thinks he sees his foes approaching, the latter, that he sees the object of his desire; and the more deeply one is under the influence of the emotion, the less similarity is required to give rise to these impressions." The passage establishes that affective residue primes formal recognition independently of the formal content. The coward and the courageous person, confronting the same minimal formal resemblance, respond differently --- not because they have access to different formal traces but because their persisting affective states modulate the threshold at which formal recognition occurs. If the affective dimension were merely a shadow of the formal trace, no such differential priming would be possible; the coward and the courageous would recognize the same minimal resemblance or fail to recognize it identically. That they respond differently is evidence that the affective residue is doing its own operative work.

\textit{De Motu Animalium} 7, 702a17-19 provides the most architecturally explicit support: "the organic parts are suitably prepared by the affections, these again by desire, and desire by imagination." Aristotle inserts the affections (\textit{pathē}) as a distinct causal step between desire and the organic parts. If the affective dimension were simply identical with the formal content of the \textit{phantasma}, its insertion as a distinct causal step would be otiose. Aristotle includes the \textit{pathē} because they do something the formal trace cannot do alone: they prepare the body for action by activating the hedonic-orectic axis. The \textit{pathē} at this point in the causal chain are the resonant \textit{orexis} at work --- the affective residue that mediates between the phantasmatic representation and the bodily action.

\textit{Rhetoric} I.11, 1370a27-b1 supplies complementary evidence in a different register. "Pleasure is the consciousness through the senses of a certain kind of emotion; but imagination is a feeble sort of sensation, and there will always be in the mind of a man who remembers or expects something the imagination of what he remembers or expects. If this is so, it is clear that memory and expectation also, being accompanied by sensation, may be accompanied by pleasure." The pleasure that accompanies memory and expectation is not a fresh evaluation performed on a retrieved formal trace; it is a residual affection carried by \textit{phantasia} as part of the image's content. The image delivers the pleasure directly --- because the pleasure was deposited alongside the formal trace in the original perceptual event, and because the dual-trace structure ensures that both persist together when the residue is activated.

The convergence of the two arguments yields the dual-trace thesis in its final form. What persists after a perceptual event is not a single undifferentiated impression but a composite \textit{pathos} with two formally distinct aspects: the resonant \textit{aisthēma} (the "what it is" of the perceived object, the sensible form preserved) and the resonant \textit{orexis} (the "how it matters for me" of the object, the hedonic-orectic valence preserved). These two aspects are one in substrate --- a single enmattered affection in the organ --- but two in being, corresponding to the distinct ways of being of the perceptual and appetitive faculties that were already numerically one in the original act. Aristotle does not name the second trace with a technical term. He names the first as \textit{aisthēma} and treats affections as \textit{pathē} generally. The dual structure as developed here is a structural requirement of his own commitments --- the commitment that sensation analytically entails affective valence (\textit{De Anima} II.3, 414b1-6), the commitment that perception and appetite are one in substrate and two in being (\textit{De Anima} III.7, 431a8-14), the commitment that the full \textit{pathos} of the sensory encounter persists after the object withdraws (\textit{On Dreams} 2, 459a24-28 and 460b2-3), and the commitment that the affective dimension has its own operative force within the causal chain of animal movement (\textit{De Motu Animalium} 7, 702a17-19). Take these commitments together, and the dual trace follows.

\subsection*{8. Synthesis: $A_2$ as Ground of $A_3$}

$A_2$, then, is the completed perceptual actuality: the joint actualization of sensible object and sense faculty in which the organ receives the form of its object without its matter, together with the resonant \textit{kinēsis} --- the dual-trace composite of resonant \textit{aisthēma} and resonant \textit{orexis} --- that the event deposits in the sensory apparatus. The section has traced $A_1$$\to$$A_2$ through six moves: the preconditions at $A_1$ and the medium that links them; the hylomorphic analysis of form-reception; the unity of mover and moved as a single actuality under two \textit{logoi}; perception as preservative \textit{alloiōsis}; the persistence of residual motion as resonant \textit{kinēsis}; and the dual-trace thesis that specifies the internal structure of the residue.

$A_2$'s architectural role follows from the iterated chain structure established at $A_0$. Each completed actuality in the chain is simultaneously the terminus of its antecedent motion and the unmoved originator of its successor. $A_2$ is the terminus of the perceptual motion from $A_1$; it is also the unmoved originator of the phantastic motion $M_2$$\to$$A_3$ through which the \textit{phantasma} proper will emerge. The resonant \textit{kinēsis} deposited at $A_2$ is the material upon which \textit{phantasia} will operate in the next stage. The \textit{phantasma} is not created ex nihilo from the sensible object's form; it is generated from the dual trace that $A_2$ has left in the sensory apparatus. \textit{Phantasia}, as Aristotle defines it at \textit{De Anima} III.3, 428b10-17, is a \textit{kinēsis} that arises from actual sensation and is necessarily similar in character to the sensation from which it arose. The similarity is grounded in the resonant \textit{kinēsis}: the \textit{phantasma} inherits both the formal aspect (resonant \textit{aisthēma}, making the image recognizable as of that object) and the affective aspect (resonant \textit{orexis}, making the image matter to the perceiver) from the residual motion deposited at $A_2$.

The generalization established at $A_0$ --- each transition from one actuality to the next is itself a motion, a further actualization of a capacity that was latent in the preceding stage --- applies here as it will at every subsequent stage. $A_2$ $\to$ $A_3$ is the motion through which the residual trace is taken up by \textit{phantasia} and crystallized into the \textit{phantasma} proper. The next section takes up that motion in detail, analysing how the resonant \textit{kinēsis} is activated by \textit{phantasia} and presented to the central organ, how the \textit{phantasma} that emerges is ontologically distinct from the trace on which it depends, and how the dual-aspect structure of the resonant \textit{kinēsis} is inherited and elaborated by the \textit{phantasma} that $A_3$ names. For now, the reader is deposited at the threshold of $A_3$, where the material deposited by perception becomes the object of the phantastic motion.\footnote{The Heideggerian analysis of perception --- particularly Heidegger's reading of \textit{aisthēsis} as a mode of Dasein's disclosive being-in-the-world and his treatment of the sensible object's presence (\textit{Anwesenheit}) as the ontological condition for perceptual encounter --- illuminates what this section establishes. Heidegger's analysis of \textit{Befindlichkeit} (state-of-mind, attunement) as the existential structure through which Dasein is always already disposed toward its world offers an especially suggestive parallel to the dual-trace thesis developed in §7, insofar as \textit{Befindlichkeit} names the affective-orectic dimension of world-disclosure that cannot be reduced to the formal cognitive aspect. A full integration of the Aristotelian and Heideggerian accounts of perception lies beyond the scope of the present chapter and is developed in [Chapter X: The Motion-Inducing Soul --- Rhetoric, Attunement, and Actualization (Heidegger, Rickert, Burke)].

\subsection*{Development Note}

Three items flag future development. First, the dual-trace thesis advanced in §7 treats the inheritance of affective valence as analytically entailed by sensation in all animals. This reading draws on \textit{De Anima} II.3, 414b1-6, which makes the entailment explicit for sensation as such. Whether every sensory affection in every animal inherits both traces, however, bears on the wider Aristotelian claim that appetite is the single source of all animal movement (\textit{De Anima} III.10, 433a27-31); and the question whether primitive organisms with only the sense of touch carry both traces or only the formal one merits engagement in a more extended treatment of animal locomotion. Nussbaum's commentary on the \textit{De Motu Animalium} and Caston's analysis of \textit{phantasia}'s role in animal cognition provide the principal points of reference for this debate.

Second, the precise mechanism by which the medium's sequential affection translates into the unified perceptual act at $A_2$ is not fully articulated in \textit{De Anima} and requires supplementary analysis from \textit{Sense and Sensibilia} and the \textit{De Motu Animalium}. The sequential character of the medium's affection (\textit{Sense and Sensibilia} 6, 447a3-7) implies that the sensible form reaches the sense-organ over a temporal extension that varies by sensory modality, yet the perceptual act itself, at the moment of completion, presents the object as a unified formal whole rather than as a sequence of partial impressions. The mechanism through which the temporally-extended transmission issues in the temporally-unified perceptual experience is a genuine aporia that the Aristotelian corpus does not fully resolve; its treatment will draw on the \textit{koinē aisthēsis} analysis at \textit{De Anima} III.1-2 and on the further articulation of \textit{phantasia}'s synthesizing role in subsequent sections.

Third, the dual-trace thesis developed here distinguishes the basic affective valence of resonant \textit{orexis} --- the pursue-or-avoid, the pleasant-or-painful --- from the determinate higher-order \textit{pathē} that Aristotle catalogues in the \textit{Rhetoric} (anger as desire accompanied by pain for conspicuous revenge; fear as disturbance at an anticipated destructive evil; pity as pain at an apparent evil befalling one who does not deserve it). The basic valence is co-constitutive with perception from $A_2$ onward and is inherited as a dispositional property by the \textit{phantasma} at $A_3$. The higher-order \textit{pathē}, by contrast, arise only downstream of \textit{doxa}'s propositional ratification at $M_2$$\to$$M_3$ --- they require the rational cognitive apparatus and the committal dimension of opinion that this section's animals (which may include non-rational animals) do not necessarily possess. The distinction between \textit{pathos} simpliciter (basic valence, resonant \textit{orexis}) and \textit{pathē} in the higher-order sense (fear, anger, pity, shame) is structurally load-bearing for the downstream sections on cognitive actuality and emotion, where it will be developed in full.}



\end{document}
